\section{来源审计：为什么 AWS 之上仍需要一层平台}

这场访谈从一个看似朴素的问题开始：既然 AWS 已经提供计算、存储、网络和数据库，为什么还需要 Vercel？回答不能停在“开发体验更好”。真正的系统问题是：\textbf{hyperscaler 暴露的是通用 primitives，而应用团队需要的是从产品意图到可运行系统的确定性转换。} 当 primitives 数量越来越多，选择、连接、配置、扩缩容和观测本身就形成新的复杂性层。

历史公开视频在 2026 年 8 月 11 日核验时已转为 private；仓库保留官方封面和 931 条时间戳字幕。字幕自动识别把 Vercel、v0、Guillermo 和若干技术名词写错，因此本讲只在上下文和官方文档一致时规范化名称。Black Friday、收入、客户损失和 uptime 等数字视为讲者现场案例，不当作审计数据。

\begin{importantbox}{本讲的两个编译/反馈闭环}
\begin{enumerate}
  \item \textbf{设计闭环}：application intent $\rightarrow$ framework semantics $\rightarrow$ platform IR $\rightarrow$ infrastructure；
  \item \textbf{运行闭环}：traffic $\rightarrow$ telemetry/metering $\rightarrow$ developer decision $\rightarrow$ new immutable deployment。
\end{enumerate}
Vercel 的价值来自同时缩短这两个闭环，而不是隐藏所有技术细节。
\end{importantbox}

先看平台在整个栈中的位置。底层 hyperscaler 解决大规模物理运营，上层框架表达应用结构，平台负责把二者连接成可持续运行的部署系统。

\lecturefigure{01-abstraction-layers.png}{平台层把应用意图编译成对 hyperscaler primitives 的受管使用。}{本地字幕 00:00--03:10；\texttt{lecture05-diagrams.py} 可复现重绘。}

\begin{knowledgebox}{读图：抽象层不是简单“包一层”}
有价值的抽象必须压缩决策空间，并提供底层 primitives 做不到的全局不变量，例如每个 commit 对应不可变部署、域名指针可原子切换、框架路由自动变成 edge routing、缓存策略与代码语义一致。若平台只是把 AWS 控制台换皮，它没有创造新的系统能力。
\end{knowledgebox}

\begin{knowledgebox}{术语消化：四个容易混淆的层次}
\begin{tabularx}{\textwidth}{L{0.20\textwidth}>{\raggedright\arraybackslash}X>{\raggedright\arraybackslash}X}
\toprule
层次 & 提供什么 & 不自动提供什么 \\
\midrule
Hyperscaler & 区域、硬件、网络和通用云服务 & 针对某个框架的完整应用语义 \\
Cloud primitive & 对象存储、函数、队列、数据库等能力 & primitives 之间的正确组合与产品默认值 \\
Platform & 受管 build、deploy、route、scale、observe & 应用业务逻辑与正确产品目标 \\
Framework & 路由、rendering、data lifecycle 等结构 & 跨区域物理基础设施和长期运营 \\
\bottomrule
\end{tabularx}
\end{knowledgebox}

\subsection{本章小结}

平台层存在的理由不是底层云“不够强”，而是通用 primitives 与应用意图之间仍有巨大翻译成本。高价值平台把这段翻译固化为可重复的编译和运营闭环。

